% ============================================================
% Introducción — traducción del original "Introduction" (1889)
% ============================================================
\chapter*{Introducción}
\addcontentsline{toc}{chapter}{Introducción}
\markboth{Introducción}{Introducción}

De todas las teorías físicas, la menos imperfecta es la de la luz, tal
como ha salido de los trabajos de Fresnel y de sus sucesores. Está
fundada, como se sabe, en la hipótesis de las ondulaciones del éter. Esta
teoría explica casi todos los hechos actualmente conocidos en óptica, y si
hay algunos que escapan a una explicación inmediata, basta con algunas
modificaciones de detalle en las hipótesis de Fresnel para dar cuenta de
ellos.

A los escépticos que piensan que algún día la teoría de las ondulaciones
correrá la misma suerte que la teoría de la emisión, se les puede
responder que, cuando Biot, en 1813, intentaba explicar los fenómenos
entonces conocidos en óptica por medio de la teoría de la emisión, no lo
lograba sino a costa de esfuerzos demasiado ingeniosos para satisfacer
plenamente al espíritu. En la teoría de la emisión había tantas hipótesis
como hechos por explicar; en la de las ondulaciones hay, es verdad, cierto
número de hipótesis, pero muchas menos que hechos explicados.

Es, pues, probable que, cualquiera que sea la suerte reservada a la
teoría de Fresnel, la mayoría de sus resultados subsistirán siempre, y que
su estudio seguirá siendo siempre útil.

En estos últimos tiempos se ha intentado sustituir la teoría de las
ondulaciones por una teoría electromagnética, que atribuye los fenómenos
ópticos a las variaciones periódicas y extremadamente rápidas que sufre un
campo magnético. Pero esta teoría conduce a los mismos resultados
analíticos que la teoría de las ondulaciones de Fresnel; solo difiere la
interpretación física de las fórmulas. Ni siquiera es imposible que estas
dos teorías terminen por concordar y no formen más que una sola.

\noindent\textbf{Objeto del curso.} — No estudiaremos sino la teoría de
las ondulaciones, siendo la teoría electromagnética inseparable de un
curso de electricidad. Nos atendremos a los trabajos de Fresnel, Cauchy,
Lamé, Briot y del señor Sarrau, en Francia; de Neumann, en Alemania; y a
los de Mac-Cullagh, en Inglaterra. Admitiremos como conocidos todos los
hechos experimentales de la óptica, y solo nos dedicaremos al desarrollo
de las teorías matemáticas.

Comenzaremos por la teoría de los pequeños movimientos en un medio
elástico; estableceremos las leyes generales del movimiento vibratorio y
de la propagación de las ondas planas; abordaremos sucesivamente el
estudio de la difracción, de las diversas teorías de la dispersión, de las
de la doble refracción, así como de la reflexión y la refracción en la
superficie de los cuerpos transparentes e isótropos, luego de los cuerpos
cristalizados, y finalmente de las superficies metálicas. Terminaremos con
el estudio de la aberración y de la propagación de la luz en un medio en
movimiento.
